% Данный файл распространяется под лицензией CC BY 4.0. Текст лицензии размещён на https://creativecommons.org/licenses/by/4.0/
% (c) Кафедра системного программирования, 2026

% Компилировать XeTeX

\documentclass[a4paper]{article}

\usepackage[a4paper, top=8mm, bottom=8mm, left=8mm, right=8mm]{geometry}

\usepackage{polyglossia}
\setdefaultlanguage[babelshorthands=true]{russian}
\setotherlanguage{english}

\usepackage{fontspec}
\setmainfont{FreeSerif}
\newfontfamily{\russianfonttt}[Scale=0.7]{DejaVuSansMono}

\usepackage[tiny, compact]{titlesec}

\usepackage{titling}
\setlength{\droptitle}{-1cm}
\pretitle{\begin{center}\begin{bfseries}\Large}
\posttitle{\par\end{bfseries}\end{center}}
\preauthor{\begin{center}\normalsize}
\postauthor{\par\end{center}\vspace{-1.8cm}}

\usepackage{hyperref}
\usepackage{bookmark}
\usepackage{csquotes}
\usepackage{booktabs}
\usepackage{array}

\title{Архиватор на базе алгоритма Хаффмана}

\author{Лиханов Даниил Викторович}

\date{}

\begin{document}

\maketitle

\begin{flushright}
    Группа: \emph{2025.Б73-мм}

    Кафедра: \emph{кафедра системного программирования}

    Научный руководитель: \emph{Литвинов Юрий Викторович}

    Номер семестра практики: \emph{2}
\end{flushright}

\section{Введение}

В современной вычислительной технике задача эффективного хранения и передачи данных
является одной из фундаментальных. Объём генерируемой информации растёт экспоненциально,
что делает алгоритмы сжатия данных без потерь критически важными инструментами для
оптимизации дискового пространства и пропускной способности каналов передачи данных.

\section{Постановка задачи}

Целью работы является разработка консольного приложения, использующего алгоритм Хаффмана
для уменьшения объёма хранимой информации без потери данных.

Актуальность работы обусловлена необходимостью эффективного хранения и передачи
информации. Сжатие данных позволяет сократить объём файлов и ускорить передачу данных по
сети. Алгоритм Хаффмана является одним из наиболее известных методов кодирования и широко
применяется в различных форматах хранения и передачи данных.

Для достижения цели требуется решить следующие задачи:

\begin{enumerate}
    \item \textbf{Выполнить обзор существующих решений:} Изучить теоретические основы алгоритма
          Хаффмана, проанализировать существующие подходы к реализации дерева Хаффмана.

    \item \textbf{Спроектировать решение:} Разработать архитектуру приложения, включая модули для
          частотного анализа входных данных, построения дерева кодирования, генерации таблицы кодов.

    \item \textbf{Реализовать решение:} Написать программный код, обеспечивающий корректное сжатие
          (архивацию) и восстановление (разархивацию) исходных файлов.

    \item \textbf{Выполнить апробацию:} Провести тестирование разработанного решения на файлах
          различного типа и размера для проверки корректности восстановления данных и оценки степени
          сжатия.
\end{enumerate}

\section{Обзор}

Сжатие без потерь (\textit{lossless compression}) применяется там, где критически важна абсолютная
точность восстановления исходных данных (тексты, программный код, исполняемые файлы). В
рамках проекта были рассмотрены три основных подхода, которые чаще всего используются в
программировании для решения этой задачи:

\subsection{Алгоритм кодирования повторов (RLE --- \textit{Run-Length Encoding})}

\textbf{Принцип работы:} Метод основан на поиске одинаковых символов или байтов, идущих подряд.
Вместо записи всей последовательности алгоритм сохраняет сам символ и количество его
повторений.

\textbf{Особенности:} Алгоритм показывает высокую эффективность только на специфических данных
(например, простых изображениях с большими одноцветными областями). На обычном тексте или
исходном коде программ данный метод неэффективен и может привести к увеличению размера
итогового файла.

\subsection{Словарный алгоритм (LZ77)}

\textbf{Принцип работы:} Алгоритм сканирует текст и ищет повторяющиеся текстовые фрагменты
(слова, фразы, сочетания букв). При обнаружении повтора вместо самого фрагмента записывается
ссылка на его предыдущее появление в тексте (смещение назад и длина строки).

\textbf{Особенности:} Обеспечивает высокий коэффициент сжатия для любых текстовых документов.
Однако его программная реализация требует создания сложных алгоритмов поиска в памяти, что
существенно усложняет структуру программы.

\subsection{Алгоритм частотного кодирования (Алгоритм Хаффмана)}

\textbf{Принцип работы:} Метод основан на подсчёте частоты появления каждого символа в файле.
Символам, которые встречаются чаще всего, присваиваются самые короткие битовые коды, а
редким символам~--- более длинные. Связь между символами и кодами строится с помощью
бинарного дерева.

\textbf{Особенности:} Показывает стабильные результаты сжатия текстовых файлов. Главная сложность
заключается в необходимости сохранять структуру дерева или таблицу частот внутри сжатого
файла, чтобы принимающая программа могла правильно восстановить данные.

Таким образом, для реализации выбран алгоритм Хаффмана как обеспечивающий стабильное
сжатие текстовых данных при относительно несложной архитектуре программы.

\section{Описание предлагаемого решения}

Программа написана на языке C++ (стандарт C++17).

Логика программы разделена на четыре последовательных шага:

\textbf{4.1 Подсчёт частот:} Чтение файла и сбор статистики появления каждого байта.

\textbf{4.2 Построение дерева:} Создание бинарного дерева Хаффмана на основе полученных частот.

\textbf{4.3 Создание таблицы кодов:} Обход дерева от корня к листьям для генерации новых битовых
кодов.

\textbf{4.4 Битовый ввод-вывод:} Упаковка данных в бинарный файл архива с помощью специального
побитового буфера.

Наибольшую сложность в проекте вызвало \textbf{управление памятью при построении дерева
Хаффмана}. Алгоритм требует постоянно извлекать из очереди два узла с минимальным
весом, объединять их и создавать для них нового общего родителя.

\section{Эксперименты}

Для оценки производительности и эффективности \textbf{алгоритма Хаффмана} был написан
специальный тестовый модуль на языке C++. В качестве исходных данных использовался набор из
шести файлов, представляющих различные типы информации. Текстовые данные разного объёма:
\texttt{small.txt} и \texttt{large.txt}. Несжатые бинарные данные: \texttt{binary.bin}. Предварительно сжатые
мультимедийные форматы и архивы: \texttt{compressed.zip}, \texttt{image.jpg}, \texttt{music.mp3}.

Для получения статистически достоверных результатов замеры времени проводились программно.
Каждая операция (сжатие~--- \texttt{encodeFile} и распаковка~--- \texttt{decodeFile}) выполнялась в цикле по 10 раз
для каждого файла. Замер времени осуществлялся с помощью библиотеки \texttt{<chrono>}.

На основе массива полученных временных интервалов программа автоматически рассчитывала:
среднее время выполнения операции (математическое ожидание) и среднеквадратичное
отклонение (СКО) для оценки стабильности работы алгоритма.

\subsection{Результаты тестирования}

Ниже представлена сводная таблица результатов работы программы (усреднённые данные за 10
прогонов). Коэффициент эффективности сжатия вычислялся как процентное отношение
уменьшенного размера к оригинальному.

\begin{table}[h]
\centering
\small
\begin{tabular}{|l|r|r|r|r|r|}
\hline
\textbf{Файл} & \textbf{\shortstack{Исходный\\размер\\(байт)}} & \textbf{\shortstack{Сжатый\\размер\\(байт)}} & \textbf{\shortstack{Эффектив-\\ность\\сжатия}} & \textbf{\shortstack{Время\\сжатия,\\среднее\\$\pm$ СКО (мс)}} & \textbf{\shortstack{Время\\распаковки,\\среднее\\$\pm$ СКО (мс)}} \\
\hline
\texttt{small.txt}        & 8\,528       & 6\,376       & +25.23\% & $9.51 \pm 3.20$          & $9.56 \pm 1.01$              \\
\texttt{large.txt}        & 16\,711\,198 & 8\,456\,310  & +49.40\% & $2\,370.60 \pm 501.44$   & $8\,800.53 \pm 1\,072.79$    \\
\texttt{binary.bin}       & 312\,688     & 228\,583     & +26.90\% & $59.45 \pm 3.95$         & $229.62 \pm 7.74$            \\
\texttt{compressed.zip}   & 4\,588\,359  & 4\,560\,423  & -0.05\%  & $912.13 \pm 446.40$      & $3\,721.26 \pm 695.45$       \\
\texttt{image.jpg}        & 2\,054\,988  & 2\,056\,590  & -0.08\%  & $465.71 \pm 41.09$       & $1\,838.93 \pm 67.51$        \\
\texttt{music.mp3}        & 10\,553\,390 & 10\,508\,615 & +0.43\%  & $2\,337.82 \pm 91.03$    & $8\,776.40 \pm 978.75$       \\
\hline
\end{tabular}
\caption{Результаты тестирования архиватора на базе алгоритма Хаффмана}
\end{table}

\textbf{Выводы:} Алгоритм блестяще справляется с текстовыми данными. Например, файл \texttt{large.txt} объёмом
почти 16 МБ был сжат практически в два раза (до 8.4 МБ). При обработке форматов \texttt{.zip}, \texttt{.jpg} и
\texttt{.mp3} размер файла остался неизменным или незначительно увеличился. Эксперимент,
проведённый в автоматическом режиме на серии из 10 запусков, Алгоритм Хаффмана
показывает высокую эффективность на избыточных данных (текстах) и демонстрирует
предсказуемое поведение на несжимаемых файлах. Производительность программы полностью
соответствует архитектуре выбранного метода кодирования.

\section{Заключение}

В ходе выполнения работы получены следующие результаты:

\begin{itemize}
    \item \textbf{6.1} Спроектировано и реализовано программное решение на языке C++ для сжатия и точного
    восстановления файлов с использованием алгоритма префиксного кодирования Хаффмана
    (модули \texttt{encoder} и \texttt{decoder}).

    \item \textbf{6.2} Разработан автоматизированный модуль тестирования для оценки скорости и стабильности
    работы реализованных алгоритмов.

    \item \textbf{6.3} Проведена успешная экспериментальная проверка программы на наборе файлов различных
    форматов и объёмов.
\end{itemize}

Репозиторий с исходным кодом проекта: \url{https://github.com/liptonn405/huffman-archive}

\end{document}